\documentclass[11pt,letterpaper]{article}
\usepackage{fontspec}
\usepackage{tocloft}
\usepackage{multicol}
\usepackage{nicefrac}
\usepackage[left=1.4in,right=1.5in,top=1in,bottom=1in]{geometry}
\setmainfont[Scale=1.2,AutoFakeBold=1.5,AutoFakeSlant=0.3]{Doves Type}

\widowpenalty=10000
\clubpenalty=10000
\interlinepenalty=500

\title{Crispy Roast Potatoes}
\author{}
\date{}

\begin{document}

\maketitle
\thispagestyle{empty}

\begin{quote}
\textit{Large russet chunks are boiled in salted water with a little baking soda so the outsides fray into a starchy paste. The chunks are tossed in garlic-rosemary oil and roasted until that paste turns into a deep-brown crust, then finished with the fried garlic, rosemary, and parsley. Serves 6--8.}
\end{quote}

\section*{Ingredients}
\setlength{\columnsep}{20pt}
\begin{multicols}{2}
\noindent
    Russet potatoes \dotfill 4 lb. \\
    Water \dotfill 2 qt. \\
    Diamond Crystal kosher salt \dotfill 2 Tbsp. \\
    Baking soda \dotfill \nicefrac{1}{2} tsp. \\
    \columnbreak
    Extra-virgin olive oil \dotfill 7 Tbsp. \\
    Fresh rosemary, finely chopped \dotfill 1 Tbsp. \\
    Garlic cloves, minced \dotfill 3 \\
    Black pepper \dotfill \nicefrac{1}{2} tsp. \\
    Fresh parsley, minced \dotfill \nicefrac{1}{4} cup \\
\end{multicols}

\section*{Directions}

\noindent
Heat the oven to \textit{450°F} (\textit{400°F} convection), with a rack in the center ---
Peel the \textbf{potatoes} and cut them into 2- to 3-inch chunks. Hold them in \textit{Large Bowl~\#1} ---
Combine 2~Tbsp. \textbf{Diamond Crystal kosher salt} and \nicefrac{1}{2}~tsp. \textbf{baking soda} in \textit{Small Bowl~\#1} (boil seasonings) ---
Mince the \textbf{garlic}. Finely chop the \textbf{rosemary}. Combine both with 6~Tbsp. \textbf{olive oil} and \nicefrac{1}{4}~tsp. \textbf{black pepper} in \textit{Small Bowl~\#2} (infusion) ---
Mince the \textbf{parsley} and set it aside in \textit{Small Bowl~\#3}

\begin{enumerate}
    \item Bring 2~qt. \textbf{water} to a rolling boil in an 8-qt. pot over high heat. The water must be boiling before the \textbf{potatoes} go in. A boiling start frays the outside while the center stays intact.

    \item Boil the \textbf{potatoes} and infuse the oil in the same window:
    \begin{enumerate}
        \item Add the \textbf{salt} and \textbf{baking soda} (\textit{Small Bowl~\#1}) and the \textbf{potatoes} (\textit{Large Bowl~\#1}). Stir so every chunk sits in the water. Add a splash of \textbf{water}, plus a pinch of \textbf{salt} and \textbf{baking soda}, when the top pieces stick out of the water. Return the pot to a boil, then lower the heat to a vigorous simmer.
        \item As soon as the simmer is steady, pour the \textbf{olive oil}, \textbf{rosemary}, \textbf{garlic}, and \textbf{pepper} (\textit{Small Bowl~\#2}) into a small saucepan. Set the empty bowl aside. Warm the oil over medium heat, stirring the whole time, until the \textbf{garlic} turns pale gold and the \textbf{rosemary} smells sweet, about \textit{2--3~minutes}. Pull the pan off the heat at the first sign of gold. Dark garlic will turn the oil bitter. Strain the oil into the empty \textit{Large Bowl~\#1}. Return the \textbf{garlic} and \textbf{rosemary} solids to \textit{Small Bowl~\#2}.
        \item At about \textit{10~minutes} after the boil returns, test a chunk. A knife should meet slight resistance in the center. The edges should look soft and frayed, and the piece should hold together when lifted. Give the pot another \textit{1~minute} when the center is still hard. Keep those frayed edges on the chunk. Starch left behind in the water will not crisp.
    \end{enumerate}

    \item Drain the \textbf{potatoes} and return them to the hot pot. Set the pot on the burner with the heat off. Let them steam for \textit{2--3~minutes}, and shake the pot once. The surfaces should look matte, and the steam should thin out. Wet surfaces stay pale in the oven. Give them another \textit{1~minute} while they still glisten. They should still be hot when you toss them.

    \item Tip the \textbf{potatoes} into \textit{Large Bowl~\#1}, onto the infused oil. Add the remaining \nicefrac{1}{4}~tsp. \textbf{black pepper} and a pinch of \textbf{salt}. Toss hard, shaking and turning, until a thick mashed-potato paste coats every chunk and the edges look shaggy. That paste becomes the crust. Add up to 1~Tbsp. more \textbf{olive oil} when the paste looks dry and powdery. Spoon off any oil that pools in the bottom of the bowl. A puddle will steam the bottoms.

    \item Spread the \textbf{potatoes} in one layer on a bare half-sheet pan, with space between the chunks. Use a second pan when they touch. A crowded pan steams them. Roast at \textit{450°F} (\textit{400°F} convection) without moving them for \textit{20~minutes}. The contact points should be golden, and a thin metal spatula should slide under each piece. Leave them another \textit{5~minutes} while they still stick. Moving them early tears off the crust.

    \item Loosen every piece, turn them, and shake the pan so the loose crumbs stay with the \textbf{potatoes}. Roast for \textit{30--40~minutes} more, turning once about halfway. With two pans, swap racks when you turn them. The crust should be deep golden brown, with a few darker edges, and it should sound crisp when tapped. The center should give under a gentle press and look fluffy when you break one open. Continue in \textit{5~minute} stretches while the crust is still pale. On convection, or with two pans, start checking at \textit{30~minutes}.

    \item Tip the \textbf{potatoes} into \textit{Large Bowl~\#1}. Add the reserved \textbf{garlic} and \textbf{rosemary} (\textit{Small Bowl~\#2}) and the \textbf{parsley} (\textit{Small Bowl~\#3}). Toss. Taste. Add \textbf{salt} and \textbf{pepper} until the crust tastes seasoned. Serve at once. They are at their best in the first few minutes. The crust softens as they cool.
\end{enumerate}

\newpage

{\footnotesize
\setlength{\columnsep}{20pt}
\setlength{\multicolsep}{6pt}
\begin{multicols}{2}
\setlength{\parindent}{0pt}
\setlength{\parskip}{2pt}
\setlength{\itemsep}{0pt}
\setlength{\parsep}{0pt}

\section*{Equipment Required}
\begin{itemize}
    \item 8-qt. pot
    \item Colander
    \item 8-qt. or larger mixing bowl (\textit{Large Bowl~\#1})
    \item Small prep bowls (3)
    \item Small saucepan
    \item Fine-mesh strainer
    \item Half-sheet pan, 18 by 13 in. (a second pan when the chunks touch)
    \item Thin metal spatula
    \item Knife and cutting board
    \item Measuring cups and spoons
    \item Kitchen scale
\end{itemize}

\section*{Hints and Notes}

\subsection*{Yield}
\begin{itemize}
    \item Serves 6--8 as a side.
    \item About \nicefrac{1}{2}~lb. of raw potato per person at 8 servings, and closer to \nicefrac{2}{3}~lb. at 6.
\end{itemize}

\subsection*{Mise en Place}
\begin{itemize}
    \item Heat the oven to \textit{450°F} (\textit{400°F} convection) before you start the pot. The boil and the roast should run as one stretch.
    \item \textit{Large Bowl~\#1} --- cut \textbf{potatoes}. After they go into the pot, strain the infused oil into this same bowl, then toss the drained \textbf{potatoes} in it.
    \item \textit{Small Bowl~\#1} --- boil seasonings: 2~Tbsp. \textbf{Diamond Crystal kosher salt}, \nicefrac{1}{2}~tsp. \textbf{baking soda}.
    \item \textit{Small Bowl~\#2} --- infusion: 6~Tbsp. \textbf{olive oil}, 1~Tbsp. \textbf{rosemary}, minced \textbf{garlic}, \nicefrac{1}{4}~tsp. \textbf{black pepper}. After straining, this bowl holds the fried \textbf{garlic} and \textbf{rosemary}.
    \item \textit{Small Bowl~\#3} --- minced \textbf{parsley}.
    \item Have 1~Tbsp. extra \textbf{olive oil} within reach for the toss.
\end{itemize}

\subsection*{Ingredient Tips}
\begin{itemize}
    \item \textbf{Russets} give the crispest crust and a fluffy center. \textbf{Yukon Golds} roast darker, taste deeper, and stay a little less crisp. Use the same weight and cut.
    \item Cut chunks at least 2~inches across. A quarter of a medium potato, or a sixth or eighth of a large one, is the right size. Small pieces overcook before the crust browns.
    \item Weigh the boil \textbf{salt} at 1~oz. \textbf{Diamond Crystal} is about 2~Tbsp. \textbf{Morton kosher salt} is denser. Use about 1~Tbsp. plus 1~tsp. Morton for the same 1~oz.
    \item Keep the \textbf{baking soda} at \nicefrac{1}{2}~tsp. More of it tastes soapy and turns the outsides to mush.
    \item \textbf{Beef fat} or \textbf{duck fat} can replace the \textbf{olive oil} in the same amount. Either one makes a richer crust.
    \item Use fresh \textbf{rosemary} and \textbf{parsley}. Dried herbs scorch in the oil and taste dusty at the end.
\end{itemize}

\subsection*{Preparation Tips}
\begin{itemize}
    \item The alkaline boil breaks down the outside of the \textbf{potato} faster than the center. That rough coat is what crisps.
    \item Leave the sheet bare. Parchment keeps the paste from frying against the metal.
    \item The first \textit{20~minutes} set the bottom crust. Wait until a spatula slides under a piece before you turn the pan.
    \item Browning time varies with the potato and the pan. Trust the color and the sound of the crust more than the clock.
    \item Two pans brown faster than one full pan. Check them at the short end of the \textit{30--40~minute} range.
    \item Break one \textbf{potato} open near the end. The center should look steamy and fluffy, and the shell should shatter.
\end{itemize}

\subsection*{Make Ahead \& Storage}
\begin{itemize}
    \item Cook from the boil through the roast in one stretch. The paste browns less evenly after it sits.
    \item Serve within a few minutes. The crust softens as the steam from the center moves outward.
    \item Refrigerate leftovers up to \textit{3~days}.
    \item Reheat uncovered on a sheet pan at \textit{450°F} for \textit{10--15~minutes}, until the crust crackles again. They will be good, and they will be less crisp than the first service.
\end{itemize}

\subsection*{Serving Suggestions}
\begin{itemize}
    \item Serve straight from the pan, beside roast meat, a fried egg, or a sharp salad.
    \item The peak is the first few minutes out of the oven.
    \item A squeeze of lemon or a pinch of chile flake can go on at the table. Add either one after the \textbf{parsley} toss.
\end{itemize}

\end{multicols}
}

\end{document}
